\chapter{Fluctuated product triple и минимальный BV-комплекс}

Gauge--KK specification gap требует не нового назначения башен, а явного
локального lift уже построенного finite-блока. Минимальный
almost-commutative кандидат в интерпретации \(K\) как четырёхмерной базы
задаётся данными
\begin{align}
 \mathcal A
 &=C^\infty(K)\otimes A_F,
 &
 A_F&=\mathbb C\oplus\mathbb C,
 \\
 \mathcal H
 &=L^2(K,S)\otimes H_8,
 &
 \mathcal D_\chi
 &=D_K\otimes\mathbf1+\gamma_K\otimes\chi D_F^*,
\end{align}
где \(H_8\) --- KO6-удвоение минимального reality-квадрата, а \(D_F^*\)
--- выбранная flattening-вакуумная орбита единичного масштаба.

\section{Gauge quotient из бимодульных зарядов}

Для сектора \(H_{ij}\) унитарий
\begin{equation}
 u=(e^{i\alpha_X},e^{i\alpha_C})\in U(A_F)
\end{equation}
действует фазой \(e^{i(\alpha_i-\alpha_j)}\). В порядке
\begin{equation}
 (H_{XX},H_{XC},H_{CX},H_{CC})
\end{equation}
вектор зарядов относительно
\(\theta=\alpha_X-\alpha_C\) равен
\begin{equation}
 q=(0,+1,-1,0).
\end{equation}
Диагональная подгруппа \(\alpha_X=\alpha_C\) действует тривиально.
Следовательно, faithful gauge quotient конечного квадрата равен
\begin{equation}
 G_F
 =\frac{U(1)_X\times U(1)_C}{U(1)_{\mathrm{diag}}}
 \simeq U(1)_{\mathrm{rel}}.
\end{equation}

Для параметров finite Dirac-блока это даёт
\begin{equation}
 x\longmapsto e^{-i\theta}x,
 \qquad
 z\longmapsto e^{+i\theta}z,
\end{equation}
поэтому произведение \(xz\), его фаза и flattening-потенциал
gauge-инвариантны.

\begin{theorem}[Relative-\(U(1)\) quotient pass]
Минимальный reality-квадрат не оставляет выбор между spectator и
динамической относительной фазой после локального lift: faithful
непрерывный quotient единственен и равен \(U(1)_{\mathrm{rel}}\).
Диагональный \(U(1)\) не связан с finite Dirac-блоком.
\end{theorem}

\section{Ковариантная производная и gauge mass}

Положим
\begin{equation}
 D_\mu x=(\partial_\mu-iA_\mu)x,
 \qquad
 D_\mu z=(\partial_\mu+iA_\mu)z.
\end{equation}
В flattening-вакууме
\begin{equation}
 |x|=|z|=\frac{\chi}{\sqrt2}
\end{equation}
trace-кинетический член даёт вдоль gauge-орбиты
\begin{equation}
 \Tr_{\mathrm{orb}}(D_\mu\Phi D^\mu\Phi)
 \supset4\chi^2 A_\mu A^\mu.
\end{equation}
Если gauge kinetic term записан как
\begin{equation}
 \frac1{4g^2}F_{\mu\nu}F^{\mu\nu},
\end{equation}
то для канонического поля \(A_\mu^{\mathrm{can}}=A_\mu/g\)
\begin{equation}
 \frac{m_A^2}{\chi^2}=c_A=8g^2.
\end{equation}
Таким образом, прежний неопределённый параметр \(c_A\) сводится к одному
каноническому gauge coupling, а не к произвольной массе.

\section{Минимальный BRST/BV complex}

Для абелева quotient достаточно полей
\begin{equation}
 (A_\mu,x,z;c,\bar c,b).
\end{equation}
BRST-дифференциал фиксируется зарядами:
\begin{align}
 sA_\mu&=\partial_\mu c,
 & sc&=0,
 \\
 sx&=-icx,
 & sz&=+icz,
 \\
 s\bar c&=b,
 & sb&=0.
\end{align}
В \(R_\xi\)-калибровке Goldstone и ghost получают одинаковый оператор
\begin{equation}
 \Delta_0+\xi m_A^2.
\end{equation}
Векторный principal block имеет вид
\begin{equation}
 \Delta_1+m_A^2
\end{equation}
с обычной \(\xi\)-зависимой продольной частью. В Feynman gauge
\(\xi=1\) четыре vector-компоненты, одна Goldstone-мода и комплексная
ghost-пара дают shorthand
\begin{equation}
 4+1-2=3
\end{equation}
для трёх физических massive-vector степеней свободы. При общем \(\xi\)
longitudinal vector, Goldstone и complex ghost pair имеют общий operator
\(\Delta_0+\xi m_A^2\) и сокращаются как
\(1+1-2=0\), оставляя те же три transverse polarizations массы \(m_A\).

\begin{proposition}[BV degree-of-freedom pass]
Gauge, Goldstone и ghost Hessian-блоки теперь происходят из одного
relative-\(U(1)\) complex. Их zero-mode supertrace вклад равен
\begin{equation}
 \Delta N_{\mathrm{gauge+ghost}}
 =3(8g^2)^2=192g^4.
\end{equation}
\end{proposition}

\section{Field-specific product operators}

Минимальный lift задаёт типы операторов:
\begin{center}
\begin{tabular}{lll}
\toprule
Сектор & Оператор на \(K\) & Bundle-данные \\
\midrule
radial/phase scalars & \(\Delta_0+m_s^2\) & trivial bosonic line \\
relative vector & \(\Delta_1+m_A^2\) & trivial adjoint line \\
Goldstone/ghost & \(\Delta_0+\xi m_A^2\) & trivial adjoint line \\
finite fermions & \(D_K+\chi D_F^*\) & spin bundle и flat character \\
\bottomrule
\end{tabular}
\end{center}
Тем самым scalar, vector и ghost спектральные типы больше не являются
свободным меню. Для абелева adjoint bundle flat relative holonomy не
скручивает gauge и ghost operators; fermion block, напротив, зависит от
spin-структуры и представления flat character.

\section{Оставшиеся нормировочные данные}

Product/BV lift закрывает gauge quotient, зарядовую единицу и структуру
Hessian-комплекса. Но он ещё не фиксирует:
\begin{enumerate}
 \item численное значение \(g\), зависящее от общей spectral
 normalization;
 \item одну из допустимых spin-структур на \(K\);
 \item flat character фермионного бимодуля;
 \item связь radion scale с \(\chi\) внутри полного quantum vacuum;
 \item общую finite-part prescription для разных Laplace-type operators.
\end{enumerate}
Поэтому полный ledger принимает более узкий вид
\begin{equation}
 B_{\mathrm{full}}
 =\frac{40+192g^4+c_\sigma^2+\Delta N_{\mathrm{KK}}
 (g,\mathfrak s,\rho_{\mathrm{flat}})}
 {64\pi^2}.
\end{equation}

\begin{theorem}[Local product-lift milestone]
По сравнению с предыдущим гейтом устранены две неопределённости:
физический gauge quotient и gauge--ghost Hessian complex. Неопределённая
gauge mass заменена вычисленной зависимостью \(c_A=8g^2\). Полное
замыкание всё ещё требует вывести \(g\), spin/flat branch и одну общую
регуляризацию KK-supertrace.
\end{theorem}

\begin{protocol}[Следующий гейт]
Следует вычислить gauge kinetic coefficient из того же \(a_4\), который
дал \(\kappa=2\), с явным orbit trace зарядов
\begin{equation}
 \Tr_{\mathrm{orb}}Q^2=2.
\end{equation}
Это должно либо зафиксировать \(g\) через один общий spectral moment, либо
доказать, что \(f_0\) остаётся новым непрерывным входом. После этого
разрешён spin/flat-branch KK audit.
\end{protocol}

Машинный аудит сохранён в
\path{s2t_v3_fluctuated_product_bv_complex_results.json}.